\subsection{扩展实验的验证与搜索协议}\label{sec:new-protocol}
九种预设方法包括常数基线、线性模型、较强 L2 正则模型、稀疏线性模型、完整决策树、随机森林、Extra Trees、直方图梯度提升和 K 近邻。回归的三种线性模型为 OLS、Ridge 和 Elastic Net；分类对应 L2 逻辑回归、较强 L2 逻辑回归与 L1 逻辑回归。森林和 Extra Trees 使用 64 棵树、深度 12、叶节点最少 5 条；提升树使用 100 次迭代并关闭内部随机早停；近邻数为 15。完整参数随源代码与拟合对象保存。

每个任务使用三折外层验证，外层训练部分再分三折搜索参数。房价和学生随机分折，种子为 42；Adult 在原训练文件内按全部非目标原始字段指纹分组并按标签分层；NBA 按完整赛季向前滚动，验证赛季晚于训练赛季。NBA 三个外层训练/验证规模分别为 3,972/2,706、6,678/2,883、9,561/3,283；最早的训练赛季不作为验证记录。其余任务每条合格记录恰好进入一次外层验证。嵌套验证将选参数与评价参数选择流程分开\citep{nestedcv}。

数值插补与标准化、分类缺失值处理和独热编码都在每次拟合的训练折完成。各方法使用相同字段、相同样本行号及相同内外层划分；房价保持美元原尺度，学生不输入同期其他考试成绩。原始记录与异常值的保留规则沿用数据核查。

网格搜索与基于 TPE 的贝叶斯超参数搜索，对同一个离散参数域各评估十二次。森林搜索最大深度 4、8 或不限，叶节点最少 3 或 10 条，候选特征比例 0.7 或 1；提升树搜索叶节点上限 7、15 或 31，L2 惩罚 0 或 1，学习率 0.05 或 0.1。两者均形成十二种组合。TPE（树结构 Parzen 估计）前四次为启动采样，后续根据已有参数及分数更新采样概率并选择参数，种子为 42 加外层折编号\citep{optunadocs}。TPE 可以重复组合，全部评估保留。分类按内层平均 AUC 选参数，回归按平均 RMSE 选参数。阈值固定为 0.5，未作阈值搜索。总计四个任务、两类模型、两种搜索、三折外层与三折内层，主搜索完成 1,728 次内层拟合；消融在此之外重复网格搜索。

\subsection{解释、消融与结果核对}
解释模型预先固定为网格调参提升树。Tree SHAP 使用训练树路径计数作为背景分布，解释全部外层验证样本，分类输出为原始对数几率\citep{shapdocs}。每条记录检查基准值加全部字段贡献能否重建原始预测。类别字段先将对应编码的有符号贡献相加，再汇总绝对重要性。

房价删去社区、面积组、质量和建成年份，并比较仅面积与质量；Adult 删去教育、工作组、婚姻和年龄；NBA 删去年龄、身体属性组和选秀年份；学生逐一删去四个背景字段。每个变体在同一外层训练记录中重新完成十二点网格搜索。消融以同折差值评价；缺少某个字段后误差变化小，也可能是其他字段提供了替代信息。SHAP 与消融用于描述预测模型，解释因果关系还需要额外的识别设计。

批量实验由独立入口运行，不调用语言模型。每折保存样本行号、内层划分、搜索轨迹、参数、拟合对象、逐条预测与 SHAP 值。另用核对脚本从逐条预测复算分类或回归指标，检查时间和分组边界、解释加和以及候选评估次数。搜索使用两个数据集进程并行，每个模型仅用一个工作线程。旧数据已参与前期分析，结果按重新定义协议的交叉验证估计报告。
